\documentclass[11pt,a4paper]{article}
% Simplified version - maximum compatibility
\usepackage[margin=1in]{geometry}
\usepackage{amsmath}
\usepackage{amssymb}
\usepackage{graphicx}
\usepackage{xcolor}
\usepackage{listings}
\usepackage{hyperref}
\usepackage{float}
\usepackage{algorithm}
\usepackage{algpseudocode}
\usepackage{setspace}

% Code highlighting
\lstset{
    basicstyle=\ttfamily\small,
    breaklines=true,
    keywordstyle=\color{blue},
    commentstyle=\color{gray},
    stringstyle=\color{red},
    showstringspaces=false,
    language=Python
}

\title{\Large \textbf{BERT Model Architecture: Mathematical Foundations and Implementation}}
\author{BERT-pytorch Protein Language Model \\ November 2025 Release}
\date{\today}

\onehalfspacing

\begin{document}

\maketitle

\begin{abstract}
This document provides a comprehensive mathematical description of the BERT (Bidirectional Encoder Representations from Transformers) model architecture, with particular emphasis on positional encoding mechanisms. We detail the mathematical formulations underlying multi-head attention, transformer blocks, and the pre-training tasks (Masked Language Model and Next Sentence Prediction). The document is tailored to the implementation in BERT-pytorch for protein language modeling.
\end{abstract}

\tableofcontents
\newpage

\section{Introduction}

BERT is a transformer-based bidirectional encoder introduced by Devlin et al. (2018) that has revolutionized natural language processing and can be applied to protein sequence analysis. Unlike traditional left-to-right language models, BERT uses masked language modeling to train a bidirectional representation of text. This document provides detailed mathematical foundations for understanding the BERT architecture.

\section{Transformer Architecture Overview}

\subsection{High-Level Architecture}

The BERT model consists of several key components:

\begin{enumerate}
    \item \textbf{Embedding Layer}: Combines token, positional, and segment embeddings
    \item \textbf{Transformer Encoder}: Stack of $n$ identical transformer blocks
    \item \textbf{Task-Specific Heads}: MLM and NSP heads for pre-training
\end{enumerate}

\noindent The forward pass through BERT can be expressed as:

\begin{equation}
\mathbf{H}_0 = \text{Embedding}(\mathbf{X}, \mathbf{S})
\end{equation}

\begin{equation}
\mathbf{H}_{i} = \text{TransformerBlock}_i(\mathbf{H}_{i-1}) \quad \text{for } i = 1, 2, \ldots, L
\end{equation}

where $\mathbf{X}$ is the input sequence, $\mathbf{S}$ is the segment information, $\mathbf{H}_i$ is the hidden representation at layer $i$, and $L$ is the number of transformer blocks.

\section{Embedding Layer}

The embedding layer transforms discrete token indices into continuous vector representations by combining three types of embeddings:

\subsection{Token Embedding}

Token embedding maps vocabulary indices to learned dense vectors:

\begin{equation}
\mathbf{e}_t = \mathbf{E}_{\text{token}}[x_t]
\end{equation}

where $\mathbf{E}_{\text{token}} \in \mathbb{R}^{|V| \times d_{\text{model}}}$ is the token embedding matrix, $|V|$ is the vocabulary size, $d_{\text{model}}$ is the hidden dimension (e.g., 256 for protein sequences), and $x_t$ is the token index at position $t$.

\subsection{Positional Embedding}

\subsubsection{Mathematical Formulation}

The positional embedding encodes absolute position information using sinusoidal functions. For each position $\text{pos}$ and dimension $i$:

\begin{equation}
PE_{(\text{pos}, 2i)} = \sin\left(\frac{\text{pos}}{10000^{2i/d_{\text{model}}}}\right)
\end{equation}

\begin{equation}
PE_{(\text{pos}, 2i+1)} = \cos\left(\frac{\text{pos}}{10000^{2i/d_{\text{model}}}}\right)
\end{equation}

where $i \in \{0, 1, \ldots, \lfloor d_{\text{model}}/2 \rfloor - 1\}$.

\subsubsection{Why Sinusoidal Positional Encoding?}

The sinusoidal positional encoding was chosen for several important reasons:

\paragraph{1. Invariance to Sequence Length}
Unlike learnable embeddings, sinusoidal encodings allow the model to extrapolate to sequences longer than those seen during training. The periodic nature of sine and cosine functions ensures smooth representation for unseen positions.

\paragraph{2. Distance Representation}
The encoding naturally captures relative positions. The dot product between positions $\text{pos}_1$ and $\text{pos}_2$ depends on their distance $\Delta = |\text{pos}_1 - \text{pos}_2|$:

\begin{equation}
PE_{\text{pos}_1} \cdot PE_{\text{pos}_2} \approx f(\Delta)
\end{equation}

This allows the attention mechanism to easily learn relative position relationships.

\paragraph{3. Unique Position Representation}
Each position has a unique encoding. We can express the positional encoding as a function of wavelengths:

\begin{equation}
\mathbf{PE}_{\text{pos}} = \begin{pmatrix}
\sin(\omega_0 \cdot \text{pos}) \\
\cos(\omega_0 \cdot \text{pos}) \\
\sin(\omega_1 \cdot \text{pos}) \\
\cos(\omega_1 \cdot \text{pos}) \\
\vdots
\end{pmatrix}
\end{equation}

where $\omega_k = 10000^{-2k/d_{\text{model}}}$ for $k = 0, 1, \ldots, d_{\text{model}}/2 - 1$.

The wavelengths increase geometrically: $\lambda_k = 2\pi \cdot 10000^{2k/d_{\text{model}}}$, ranging from $\lambda_0 = 2\pi$ to $\lambda_{\max} = 2\pi \cdot 10000$.

\paragraph{4. Linear Transformation Property}
A key property is that the positional encoding for position $\text{pos} + \Delta$ can be expressed as a linear transformation of $PE_{\text{pos}}$:

\begin{equation}
PE_{(\text{pos}+\Delta)} = \mathbf{M}(\Delta) \cdot PE_{(\text{pos})}
\end{equation}

where $\mathbf{M}(\Delta)$ is a rotation matrix that depends only on the distance $\Delta$. This enables the model to learn relative position biases efficiently.

\subsection{Segment Embedding}

Segment embeddings distinguish between two input sequences (for tasks involving sentence pairs):

\begin{equation}
\mathbf{e}^{\text{seg}}_t = \mathbf{E}_{\text{segment}}[s_t]
\end{equation}

where $\mathbf{E}_{\text{segment}} \in \mathbb{R}^{S \times d_{\text{model}}}$ with $S=3$ possible values (pad=0, segment A=1, segment B=2).

\subsection{Combined Embedding}

The final embedding at position $t$ is the sum of all three embeddings:

\begin{equation}
\mathbf{z}_t = \mathbf{e}_t + \mathbf{PE}_t + \mathbf{e}^{\text{seg}}_t
\end{equation}

This is followed by layer normalization and dropout:

\begin{equation}
\mathbf{x}_0^{(t)} = \text{Dropout}(\text{LayerNorm}(\mathbf{z}_t))
\end{equation}

\section{Multi-Head Attention Mechanism}

\subsection{Scaled Dot-Product Attention}

The foundation of the transformer is the scaled dot-product attention mechanism:

\begin{equation}
\text{Attention}(\mathbf{Q}, \mathbf{K}, \mathbf{V}) = \text{softmax}\left(\frac{\mathbf{Q}\mathbf{K}^T}{\sqrt{d_k}}\right)\mathbf{V}
\end{equation}

where:
\begin{itemize}
    \item $\mathbf{Q} \in \mathbb{R}^{n \times d_k}$ is the query matrix
    \item $\mathbf{K} \in \mathbb{R}^{m \times d_k}$ is the key matrix
    \item $\mathbf{V} \in \mathbb{R}^{m \times d_v}$ is the value matrix
    \item $n$ is the sequence length
    \item $d_k$ is the key dimension (typically $d_k = d_v = d_{\text{model}}/h$ where $h$ is the number of heads)
\end{itemize}

\subsubsection{Scaling Factor}

The scaling factor $\frac{1}{\sqrt{d_k}}$ is crucial for numerical stability:

\begin{equation}
\frac{\mathbf{Q}\mathbf{K}^T}{\sqrt{d_k}} \in \mathbb{R}^{n \times m}
\end{equation}

Without scaling, the dot products grow large with dimensionality, pushing the softmax into regions with extremely small gradients. Scaling ensures:

\begin{equation}
\text{Var}\left(\frac{\mathbf{Q}\mathbf{K}^T}{\sqrt{d_k}}\right) \approx 1
\end{equation}

\subsection{Multi-Head Attention}

Instead of performing attention once with $d_{\text{model}}$-dimensional vectors, multi-head attention performs $h$ parallel attention operations:

\begin{equation}
\text{head}_i = \text{Attention}(\mathbf{Q}\mathbf{W}_i^Q, \mathbf{K}\mathbf{W}_i^K, \mathbf{V}\mathbf{W}_i^V)
\end{equation}

where:
\begin{itemize}
    \item $\mathbf{W}_i^Q, \mathbf{W}_i^K, \mathbf{V}_i^V \in \mathbb{R}^{d_{\text{model}} \times d_k}$ are learnable projection matrices
    \item $i = 1, 2, \ldots, h$ is the head index
    \item $h$ is the number of heads (typically 8)
\end{itemize}

The outputs are concatenated and projected:

\begin{equation}
\text{MultiHead}(\mathbf{Q}, \mathbf{K}, \mathbf{V}) = \text{Concat}(\text{head}_1, \ldots, \text{head}_h)\mathbf{W}^O
\end{equation}

where $\mathbf{W}^O \in \mathbb{R}^{d_{\text{model}} \times d_{\text{model}}}$.

\section{Transformer Block}

\subsection{Architecture}

Each transformer block consists of two sub-layers with residual connections and layer normalization:

\begin{equation}
\mathbf{z}^{(l)} = \mathbf{x}^{(l)} + \text{MultiHeadAttn}(\text{LayerNorm}(\mathbf{x}^{(l)}))
\end{equation}

\begin{equation}
\mathbf{x}^{(l+1)} = \mathbf{z}^{(l)} + \text{FFN}(\text{LayerNorm}(\mathbf{z}^{(l)}))
\end{equation}

where $l$ is the layer index and FFN is the feed-forward network.

\subsection{Feed-Forward Network}

The feed-forward network applies two linear transformations with a non-linearity:

\begin{equation}
\text{FFN}(\mathbf{x}) = \mathbf{W}_2(\sigma(\mathbf{W}_1\mathbf{x} + \mathbf{b}_1)) + \mathbf{b}_2
\end{equation}

where:
\begin{itemize}
    \item $\mathbf{W}_1 \in \mathbb{R}^{d_{\text{model}} \times d_{\text{ff}}}$ with $d_{\text{ff}} = 4 \cdot d_{\text{model}}$ (expanded dimension)
    \item $\mathbf{W}_2 \in \mathbb{R}^{d_{\text{ff}} \times d_{\text{model}}}$ (projection back)
    \item $\sigma$ is the GELU activation function
\end{itemize}

GELU is defined as:

\begin{equation}
\text{GELU}(x) = x \cdot \Phi(x)
\end{equation}

where $\Phi(x)$ is the cumulative distribution function of the standard normal distribution. The approximation used in this implementation is:

\begin{equation}
\text{GELU}(x) = 0.5x\left(1 + \tanh\left(\sqrt{\frac{2}{\pi}}(x + 0.044715x^3)\right)\right)
\end{equation}

\subsection{Layer Normalization}

Layer normalization stabilizes training by normalizing the activations across the feature dimension:

\begin{equation}
\text{LayerNorm}(\mathbf{x}) = \gamma \odot \frac{\mathbf{x} - \mu}{\sqrt{\sigma^2 + \epsilon}} + \beta
\end{equation}

where:
\begin{itemize}
    \item $\mu = \frac{1}{d_{\text{model}}}\sum_{i=1}^{d_{\text{model}}} x_i$ is the mean across features
    \item $\sigma^2 = \frac{1}{d_{\text{model}}}\sum_{i=1}^{d_{\text{model}}} (x_i - \mu)^2$ is the variance
    \item $\gamma, \beta$ are learnable affine parameters
    \item $\epsilon = 10^{-6}$ is a small constant for numerical stability
\end{itemize}

\section{Pre-training Tasks}

\subsection{Masked Language Model (MLM)}

The MLM task trains the model to predict masked tokens based on surrounding context:

\begin{equation}
\mathcal{L}_{\text{MLM}} = -\sum_{t \in M} \log P(x_t | \mathbf{x}_{\setminus t})
\end{equation}

where $M$ is the set of masked token positions.

For each sequence, 15\% of tokens are randomly selected for masking:
\begin{itemize}
    \item 80\% of the time: replaced with [MASK] token
    \item 10\% of the time: replaced with a random token
    \item 10\% of the time: kept unchanged
\end{itemize}

\subsection{Next Sentence Prediction (NSP)}

The NSP task trains the model to determine if two sentences are consecutive:

\begin{equation}
\mathcal{L}_{\text{NSP}} = -\log P(\text{IsNext} | \mathbf{h}_0)
\end{equation}

where $\mathbf{h}_0$ is the hidden representation of the [CLS] token.

For each input sequence:
\begin{itemize}
    \item 50\% of the time: two consecutive sentences (label = 1)
    \item 50\% of the time: two random sentences (label = 0)
\end{itemize}

\subsection{Joint Training Objective}

The total pre-training loss combines both tasks:

\begin{equation}
\mathcal{L}_{\text{total}} = \mathcal{L}_{\text{MLM}} + \mathcal{L}_{\text{NSP}}
\end{equation}

\section{Implementation Details}

\subsection{Key Hyperparameters}

\begin{table}[H]
\centering
\begin{tabular}{|l|c|c|}
\hline
\textbf{Parameter} & \textbf{Protein Default} & \textbf{BERT-Base} \\
\hline
Hidden Size & 256 & 768 \\
Number of Layers & 8 & 12 \\
Number of Attention Heads & 8 & 12 \\
Feed-Forward Dimension & 1024 & 3072 \\
Dropout Rate & 0.1 & 0.1 \\
Max Sequence Length & 512 & 512 \\
Vocabulary Size & 26 & 30,522 \\
\hline
\end{tabular}
\end{table}

\section{Conclusion}

The BERT architecture combines several elegant mathematical principles for effective language modeling that transfer naturally to protein sequence analysis.

\end{document}
